% Source: 27_Chapter_18_The_Search_for_Unity_Charlemagne_and_Ch.txt
\bookchapter[The Search for Unity]{CHAPTER 18}{THE SEARCH FOR UNITY}{Charlemagne and Christendom}

\leadin{T}{HE DATE WAS APRIL} 25, 799, St. Mark’s Day, a customary day of repentance and prayer.
This year, however, the faithful had been plagued by accidents and severe crop damage. So Pope Leo
III (795–816) was leading a procession through Rome for God’s blessing upon the fields and their
harvest.

The procession led from the Church of the Lateran through the middle of the city to St. Peter’s. As
it turned the corner by the monastery of Saints Stephen and Silvester, armed men rushed at the pope.
They drove off his attendants and pulled Leo off his horse and whisked him away to a Greek
monastery.

It was mutiny led by ranking officials and dignitaries loyal to the previous pope, Adrian I. Perjury
and adultery were among the charges leveled at Leo. Using the cover of darkness, the pope’s
supporters were able to rescue him and to bring him back to St. Peter’s. But continued fighting in
the streets convinced Leo that he needed outside help. His appeal went to the traditional protector
of the papacy, the king of the Franks, Charles the Great.

The next year Charles crossed the Alps with an army, prepared to settle the pope’s problem once and
for all. In December the king presided over a large assembly of bishops, nobles, diplomats, members
of the royal household, and rebels. On December 23 the pope, holding the gospel in his ringed hand,
took an oath purging himself of the accusations against him. With that the mutiny waned, but the
stage was set for a more important development.

On Christmas Day, Charles came to St. Peter’s with a large retinue for the Christmas worship. Leo
sang the mass and Charles prayed on his knees in front of the crypt of the apostle. Charles saw the
pope approach. In his hands was a golden crown. Leo placed it on Charles’s head as the congregation
cried, “To Charles, the most pious, crowned Augustus by God, to the great peace-making Emperor, long
life and victory!” The pope prostrated himself. Charles the Great, King of the Franks, had restored
the Christian Roman Empire.

Modern times are marked by the idea of autonomous, sovereign states without religious affiliation
and by the concept of the church as a voluntary association apart from the rest of organized
society. But neither of these ideas existed in the Middle Ages.

Drawing upon Augustine’s vision of the “City of God,” Charles the Great engrafted the Christian
concept of a universal Catholic Church on the stock of the traditional Roman view of empire and gave
to the medieval world Christendom, a unified society mingling religious (or eternal) concerns with
earthly (or temporal) affairs.

How did it happen? How could the kingdom that Jesus said was “not of the world” become so much a
part of worldly power?

The answer lies in the persistence of an idea and the rise of a powerful kingdom.

\section{IDEAS DIE SLOWLY}

Centuries after its fall to barbarism, the Roman Empire in the West continued its sway over the
imaginations of men. The barbarians had many kingdoms and were often at war with one another. But
men still longed for the unity that had once marked the empire, and they looked for the day when a
new Roman Empire would appear. As the Greeks believed that Rome had passed over to Constantinople,
so the Roman people and their German neighbors thought that the empire would live again among them.

In the blending of Roman and Germanic peoples and cultures, the Franks ascended above the others and
seemed destined to restore the imperial authority. Clovis, with the active support of the Catholic
Church, had made the kingdom of the Franks a dominant power among the German tribes.

After the death of Clovis, however, his dynasty began to decay from inner weakness. The Germanic
practice of treating the kingdom as personal property and dividing it among the king’s sons resulted
in constant and bitter civil war. So the royal heirs plotted murders and became adept at intrigue
and treachery.

At the same time a new center of power arose from the landed aristocrats. More and more authority
shifted into their hands. Among these powerful landowners, one emerged as the most influential
figure in the kingdom. He was called “mayor of the palace.”

A new day dawned for the kingdom of the Franks in 714, when Charles Martel, Charlemagne’s
grandfather, became mayor of the palace. Martel allowed the Merovingian kings to retain their claim
to the throne. But they were mere figureheads; the real power rested with the mayor of the palace.

Many students remember Charles for his victory over the Muslim invaders of Europe. This triumph
earned him the surname Martel, “The Hammer.” In 711 a Muslim army from North Africa had invaded
Spain, and by 718 the weak kingdom of the Visigoths had collapsed. With most of the peninsula under
their control, the Muslims began making raids across the Pyrenees Mountains. In 732 Charles Martel
met them near Tours, deep within the Frankish kingdom. He inflicted heavy losses upon them, so
during the night they retreated toward Spain and were never again a major threat to central Europe.

Charles Martel’s son, Pepin the Short (741–768), was a worthy successor to his father. He thought
the time had come, however, to legalize the regal power exercised by the mayors of the palace. He
turned to the pope for a ruling stipulating that whoever had the actual power should be the legal
ruler. He got what he wanted. With the papal blessing, Boniface, the great English missionary among
the Germans, in 751 crowned Pepin king of the Franks. The last Merovingian was quietly shelved in a
secluded monastery. Three years later the pope blessed this coup d’etat by crossing the Alps and
personally anointing Pepin, in the Old Testament manner, as the Chosen of the Lord.

As one historian explains it, behind the pope’s action lay his need for a powerful protector. In 751
the Lombards had conquered the imperial territory at Ravenna, the seat of Byzantine government in
Italy, and were demanding tribute from the pope and threatening to take Rome. Following Pepin’s
coronation, the pope secured his promise of armed intervention in Italy and his pledge to give the
papacy the territory of Ravenna once it was conquered. In 756 a Frankish army forced the Lombard
king to surrender his conquests, and Pepin officially conferred the Ravenna territory upon the pope.
Known as the “Donation of Pepin,” the gift made the pope a temporal ruler over the Papal States, a
strip of territory that extended diagonally across Italy from coast to coast. Peter recovered his
sword.

This alliance between the Franks and the papacy affected the course of European politics and
Christianity for centuries. It accelerated the separation of the Latin from the Greek Church by
providing the papacy with a dependable western ally in place of the Byzantines, hitherto its only
protector against the Lombards; it created the Papal States, which played a major role in Italian
politics until the late nineteenth century; and, by the ritual anointing, it provided western
kingship with a religious sanction that would in time contribute to the rivalry between pope and
emperor.

Only one significant step remained to restore the Christian Empire to the West. That came with the
rise to power of Pepin’s illustrious son, Charles. Today, we know him as Charlemagne, Charles the
Great.

When Charles succeeded his father in 768 his mind was set on three goals: military power to crush
his enemies, religious power to direct his people’s souls, and intellectual power to instruct both
souls and minds. Charlemagne’s success in these areas made Europe—the new political order—nominally
Christian, for better or for worse, for a thousand years.

\section{ARCHITECT OF AN EMPIRE}

\begin{blockquote}
Einhard, in his famous biography of Charlemagne, pictured his king as a natural leader of men—tall,
physically strong, and a great horseman who was always in the van of the hunt. Although he was
preeminently a successful warrior-king, leading his armies on yearly campaigns, Charlemagne also
sought to provide an effective administration for his realm.
\end{blockquote}

Four areas were successfully annexed to his kingdom by military might. First was his southern
border. Taking advantage of feuds among the Muslims in Spain, Charlemagne sought to extend
Christendom into that land. In 778 he crossed the Pyrenees with only minor success. On later
expeditions, however, he drove the Muslims back to the Ebro River and established a frontier area
known as the Spanish Mark (or March) centered around Barcelona.

Second, Charlemagne conquered the Bavarians and the Saxons, the last of the independent Germanic
tribes.

\begin{blockquote}
It took thirty-two campaigns to subdue the staunchly pagan Saxons, who lived between the Rhine and
Elbe rivers. Charlemagne divided Saxony into bishoprics, built monasteries, and proclaimed harsh
laws against paganism. Eating meat during Lent, cremating the dead (an old pagan practice), and
pretending to be baptized were offenses punishable by death.
\end{blockquote}

His brutality in conquering and converting the Saxons was extreme by any measure.

A third trouble spot was the kingdom’s eastern frontier, which was continually threatened by the
Slavs and the Avars, Asiatic nomads related to the Huns. “In six campaigns Charlemagne decimated the
Avars and then set up his own military province in the valley of the Danube to guard against any
possible future plundering by eastern nomads. Called the East Mark, this territory later became
Austria.”

Finally, like his father before him, Charlemagne intervened in Italian politics. Expansionist
ambition drove the Lombard king to invade again the papal territories. At the behest of the pope,
Charlemagne defeated the Lombards in 774 and proclaimed himself their king. While in Italy, he
cemented his father’s alliance with the Church of Rome by confirming the Donation of Pepin.

That first incursion into Italy proved to be the prelude to the one in 800 that resulted in his
coronation as emperor. The pope needed protection. Charlemagne needed divine sanction.

After his Christmas coronation Charlemagne said that he did not know it had been planned. But he
lived up to it. He wrote in his dispatches, “Charles, by the will of God, Roman Emperor,
Augustus~.~.~. in the year of our consulship I.” He had an oath taken to him as Caesar by all
officers, lay or ecclesiastical. He sent ambassadors to soothe the anger of the emperor in
Constantinople and in 812 the Eastern court acknowledged him.

The ceremony in St. Peter’s demonstrated that the memory of the Roman Empire survived as a vital
tradition in Europe and that there was a strong desire to reestablish political unity. The
coronation also inaugurated, however, a long-standing struggle between the revived empire and the
papacy.

In medieval theory church and state were but two aspects of Christendom: the one representing
Christian society organized to secure spiritual blessings, the other the same society united to
safeguard justice and human welfare. Theoretically church and state were in harmonious interplay,
each aiming to secure the good of mankind.

In fact, however, the pope and emperor were contestants. The ever-present question was, should the
church rule the state or the state control the church? This contest was illustrated on countless
fields, large and small, throughout the Middle Ages. Since time depended upon eternity, the papal
party held that the emperor depended on the pope. But since Constantine and Charlemagne revealed
that God also ordained the Christian state, the imperial party argued that the emperor was
independent of the pope; and if the Holy Father neglected his eternal office in time, the emperor
might even correct or control the pope.

In his lifetime Charlemagne left no doubts about sovereignty in his empire. He gave Europe a
dominating father figure. Everyone within the realm was answerable to him. To solve the problem of
supervising local officials, a concern that plagued all German rulers, Charlemagne issued an
ordinance creating the missi dominici, the king’s envoys. Pairs of these itinerant officials,
usually a bishop and a lay noble, traveled throughout the realm to check on the local
administration. Even the pope could not hide from Charles’s watchful eyes.

Charlemagne also fostered a revival of learning and the arts. His efforts have prompted historians
to speak of this period as a “cultural rebirth.” In 789 Charlemagne decreed that every monastery
must have a school for the education of boys in “singing, arithmetic, and grammar.” At
Aix-la-Chapelle, his capital, the emperor sponsored a palace school for the education of the royal
household and the stimulation of learning throughout the realm. His importation of scholars from
Italy and Ireland promoted studies in Greek and Latin. Alcuin, an Anglo-Saxon scholar in charge of
the school, began the arduous task of reviving learning by undertaking the first step of writing
textbooks on grammar, spelling, rhetoric, and logic. Alcuin exhorted his students, “Ye lads whose
age is fitted for reading, learn! The years go by like running water. Waste not the teachable days
in idleness!”

Few historians challenge Charlemagne’s claim to greatness as a major constructive figure of world
history. From the new center in the north, rather than on the Mediterranean, he extended Christian
civilization in Europe. After three centuries of disorder, he restored a measure of law and order.
His patronage of learning left a cultural heritage that later generations could build upon. And the
imperial ideal that he revived persisted as a political force in Europe until 1806, when the Holy
Roman Empire was terminated by another self-styled emperor, Napoleon Bonaparte.

\section{THE DECLINE INTO FEUDALISM}

Unfortunately, Charlemagne’s empire afforded no more than a breathing space. Its territories were
too vast and its nobility too powerful to be held together under existing conditions after the
dominating personality of its creator had passed from the scene. Under his weak successors the
empire disintegrated amid the confusion of civil wars and devastating new invasions. When Vikings
began sweeping out of the Northland, people increasingly surrendered both their lands and their
persons to the many counts, dukes, and other local lords in return for protection. These
disintegrating conditions presented a new challenge to the church and to the unity of Europe. We
call it feudalism.

Feudalism was a type of government in which political power was exercised locally by private
individuals rather than by the agents of a centralized state. One popular text explains it this way:
“Fully developed feudalism was a fusion of three basic elements: (1) the personal element, called
lordship or vassalage, by which one nobleman, the vassal, became the loyal follower of a stronger
nobleman, the lord (or suzerain); (2) the property element, called the fief (usually land), which
the vassal received from his lord to enable him to fulfill the obligations of vassalage; and (3) the
governmental element, meaning the private exercise of governmental functions over vassals and fiefs.
The roots of these three elements ran back to late Roman and early Germanic times.”

Central to feudalism was the personal bond between lord and vassal. In the ceremony known as the act
of homage, the vassal knelt before his lord and promised to be his “man.” In the oath of fealty that
followed, the vassal swore on the Bible, or some other sacred object, that he would remain true to
his lord. Next, in the ritual of investiture, a lance, a glove, or even a bit of straw was handed
the vassal to signify his jurisdiction (not ownership) over the fief.

\begin{sidebar}{Praying Hands}
By the end of the eleventh century, the ceremony included the vassal kneeling before the feudal
lord. The vassal would join his hands and extend them. The feudal lord would clasp his hands around
the vassal’s hands, pull him up, and kiss him. The vassal’s hands joined and extended have become
the symbol and posture of prayer today.
\end{sidebar}

The feudal contract thus entered into by lord and vassal was considered sacred and binding upon both
parties. Breaking this tie of mutual obligations was considered a felony, because it was the
fundamental bond of early medieval society. The lord, for his part, was obliged to give his vassal
protection and justice. The vassal’s primary duty was military service. He was expected to devote
forty days’ service each year to the lord without pay.

Since the church was so much a part of medieval life, it could not escape inclusion in the feudal
system. The unsettled conditions caused by new invaders (Vikings from the north and Magyars from
Asia) forced church officials to enter into close relations with the only power able to offer them
protection: the feudal barons in France and the kings in Germany. Bishops and abbots thus became
vassals, receiving fiefs for which they were obligated to provide the usual feudal services. This
loyalty to higher lords created unusual conflicts for those bishops who looked to the pope as God’s
appointed shepherd of the church. In the tenth and early eleventh centuries the pope was in no
position to challenge anyone. The office fell into decay after becoming a prize sought by local
Roman nobles.

On the positive side, however, the church in time sought to influence for the better the behavior of
the feudal barons. In addition to attempting to add Christian virtues to the code of knightly
conduct called chivalry, the church tried to impose limitations on feudal warfare. In the eleventh
century bishops inaugurated the Peace of God and Truce of God movements. The Peace of God banned
from the sacraments all persons who pillaged sacred places or refused to spare noncombatants. The
Truce of God established “closed seasons” on fighting: from sunset on Wednesday to sunrise on Monday
and certain longer periods, such as Lent. Unfortunately, both movements were generally ineffective.

Only after the German king Otto the Great revived the Roman Empire in the West in 962 was some sense
of unity restored. With the renewal of the empire, however, came the old rivalry between church and
state.

\begin{blockquote}
Otto claimed to be the successor of Augustus, Constantine, and Charlemagne, although his actual
power was confined to Germany and Italy. At first the papacy looked to the German king for
protection against the unruly Italian nobles who for a century had been making a prize of the
papacy. From the church’s viewpoint, however, this arrangement had its drawbacks, for the German
kings continued to interfere in ecclesiastical affairs—even in the election of popes.
\end{blockquote}

During the eleventh century the controversy between church and state centered on the problem of lay
investiture. Theoretically, on assuming office a bishop or abbot was subject to two investitures:
his spiritual authority was bestowed by a church official and his feudal or civil authority by the
king or a noble. In actual fact, however, feudal lords and kings came to control both the
appointment and the installation of churchmen. This practice was most pronounced in Germany, where
control of the church was the foundation of the king’s power. The German church was in essence a
state church.

\section{SETTING IN ORDER THE HOUSE OF GOD}

The church was ill prepared to challenge kings and emperors; it needed to set its own house in
spiritual order. This began with a far-reaching revival within the reformed Benedictine order of
Cluny, founded in 910. From the original monastery in Burgundy, a powerful impulse radiated for the
reform of the feudalized church. The Cluniac program began as a movement for monastic reform, but in
time it called for the enforcement of clerical celibacy and the abolition of simony, the purchase or
sale of a church office. (The term simony comes from Simon the magician, who tried to buy the gift
of the Holy Spirit from the apostles, Acts 8:9–25.) The ultimate goal of the Cluniac reformers was
to free the entire church from secular control and subject it to papal authority. Some three hundred
Cluniac houses were freed from lay control, and in 1059 the papacy itself was removed from secular
interference by the creation of the College of Cardinals, which henceforth elected the popes.

The man behind the reform of the papacy was an arch-deacon named Hildebrand. In 1073 he was elected
pope, and as Gregory VII (1073–1085) he claimed unprecedented power for the papacy. Gregory held as
his ideal the creation of a Christian commonwealth under papal control. Instead of conceding
equality between the church and the state, he insisted that the spiritual power was supreme over the
temporal. In 1075 he formally prohibited lay investiture and threatened to excommunicate any layman
who performed it and any churchman who submitted to it. This drastic act virtually declared war
against Europe’s rulers, since most of them practiced lay investiture. The climax to the struggle
occurred in Gregory’s clash with the emperor Henry IV.

The pope accused Henry of simony and lay investiture in appointing his own choice to be the
archbishop of Milan. Gregory summoned Henry to Rome to explain his conduct. Henry’s answer was to
convene in 1076 a synod of German bishops that declared Gregory a usurper and unfit to occupy the
Roman See: “Wherefore henceforth we renounce, now and for the future, all obedience unto thee.” In
retaliation Gregory excommunicated Henry and deposed him, absolving his subjects from their oaths of
allegiance.

At last, driven to make peace with the Holy Father by a revolt among the German nobles, Henry
appeared before Gregory in January 1077 at Canossa, a castle in the mountains of Italy. Dressed as a
penitent, the emperor stood barefoot in the snow for three days and begged forgiveness until, in
Gregory’s words “We loosed the chain of the anathema and at length received him~.~.~. into the lap
of the Holy Mother Church.”

“This dramatic humiliation of the emperor,” says one historian,

\begin{blockquote}
did not resolve the quarrel, nor do contemporary accounts attach much significance to the
incident—public penance was not uncommon in those days even for kings. Yet the pope had made
progress toward freeing the church from interference by laymen and toward increasing the power and
prestige of the papacy. The problem of lay investiture was settled in 1122 by the compromise known
as the Concordat of Worms. The church maintained the right to elect the holder of an ecclesiastical
office, but only in the presence of the emperor or his representative.
\end{blockquote}

Later popes added little to Gregory’s theories of the office. They, like Gregory, insisted that
Christian society was organized under the pope, its visible head, and he was guarded against all
possibility of error by the presence of Peter perpetually present in his successors, the bishops of
Rome.

Many of Gregory’s claims now appear intolerable. Yet we must agree that Gregory and his powerful
successors stood for two principles that to the Christian are incontestable: (1) In the loyalties of
men, the spiritual has the primacy over the secular; and (2) the families of men can find true unity
only in Christ and in obedience to the law of God. Medieval society was far from perfect. But during
the Middle Ages, Europe became conscious of itself as a unity, far beyond the uncertain limits of
the Holy Roman Empire; and the church attained a level of power and influence over the lives of men,
for the most part used helpfully, such as it has never known since that time.

\readinghead{Suggestions for Further Reading}
\begin{readinglist}
\entry Baldwin, Marshall W. The Mediaeval Church. Ithaca, NY: Cornell University Press, 1953.
\entry Barraclough, Geoffrey. The Crucible of Europe. Berkeley: University of California Press,
  1976.
\entry Barraclough, Geoffrey. The Medieval Papacy. New York: Harcourt, Brace \& World, 1968.
\entry Fremantle, Anne. Age of Faith. New York: Time-Life Books, 1968.
\entry *Logan, Donald. History of the Church in the Middle Ages. New York: Routledge, 2003.
\entry Noll, Mark A. Turning Points: Decisive Moments in the History of Christianity. Grand Rapids:
  Baker Books, 2000.
\entry Russell, Jeffrey Burton. A History of Medieval Christianity: Prophecy and Order. Arlington
  Heights, IL: AHM Publishing Corporation, 1968.
\entry *Southern, R. W. Western Society and the Church in the Middle Ages: The Penguin History of
  the Church, Vol. 2. New York: Penguin, 1990.
\end{readinglist}
